\SourceLeaf{145}
\SixthArticle{14}

Es wird sich lohnen, hier einige Sonderfälle anzuführen. Setzt man $t=1$, $\beta=1$ in Formel I des Artikels 5, so ergibt sich
\[
[31]\qquad
(1+u)^n=
\cfrac{1}{1-\cfrac{nu}{1+\cfrac{\frac{n+1}{2}u}{1-\cfrac{\frac{n-1}{2\cdot3}u}{1+\cfrac{\frac{2(n+2)}{3\cdot4}u}{1-\cfrac{\frac{2(n-2)}{4\cdot5}u}{1+\mathrm{usw.}}}}}}}
\]

Aus den Formeln VI, VII des Artikels 5 folgt
\[
[32]\qquad
\log(1+t)=
\cfrac{t}{1+\cfrac{\frac12t}{1+\cfrac{\frac16t}{1+\cfrac{\frac26t}{1+\cfrac{\frac2{10}t}{1+\cfrac{\frac3{10}t}{1+\cfrac{\frac3{14}t}{1+\mathrm{usw.}}}}}}}}
\]
\[
[33]\qquad
\log\cfrac{1+t}{1-t}=
\cfrac{2t}{1-\cfrac{\frac13tt}{1-\cfrac{\frac{2\cdot2}{3\cdot5}tt}{1-\cfrac{\frac{3\cdot3}{5\cdot7}tt}{1-\cfrac{\frac{4\cdot4}{7\cdot9}tt}{1-\mathrm{usw.}}}}}}
\]
Ersetzt man hier die Zeichen $-$ durch $+$, so erhält man einen Kettenbruch für den Arkustangens von $t$.
\par
Ferner haben wir
\[
[34]\qquad
e^t=
\cfrac{1}{1-\cfrac{t}{1+\cfrac{\frac13t}{1-\cfrac{\frac16t}{1+\cfrac{\frac16t}{1-\cfrac{\frac1{10}t}{1+\cfrac{\frac1{10}t}{1-\mathrm{usw.}}}}}}}}
\]
\SourceLeaf{146}
\[
[35]\qquad
t=
\cfrac{\sin t\cos t}{1-\cfrac{\frac{1\cdot2}{1\cdot3}\sin t^2}{1-\cfrac{\frac{1\cdot2}{3\cdot5}\sin t^2}{1-\cfrac{\frac{3\cdot4}{5\cdot7}\sin t^2}{1-\cfrac{\frac{3\cdot4}{7\cdot9}\sin t^2}{1-\cfrac{\frac{5\cdot6}{9\cdot11}\sin t^2}{1-\mathrm{usw.}}}}}}}
\]

Setzt man $\alpha=3$, $\gamma=\frac52$ in Formel 30, so ergibt sich unmittelbar der in Artikel 90 der \textit{Theoria motus corporum coelestium} angegebene Kettenbruch. Dort wurden zwei weitere Kettenbrüche angeführt, deren Herleitung bei dieser Gelegenheit nachzuliefern angemessen erschien. Setzt man
\[
Q=1-\cfrac{\frac{5\cdot8}{7\cdot9}x}{1-\cfrac{\frac{1\cdot4}{9\cdot11}x}{1-\frac{7\cdot10}{11\cdot13}x\ \mathrm{usw.}}}
\]
so ergibt sich an der angeführten Stelle $\displaystyle x-\xi=\cfrac{x}{1+\frac{2x}{35Q}}=\frac{xQ}{Q+\frac2{35}x}$ und daher
\[
\xi=\frac{\frac2{35}xx}{Q+\frac2{35}x}
\]
Dies ist die erste Formel; die zweite wird auf folgende Weise hergeleitet. Setzt man
\[
R=1-\cfrac{\frac{1\cdot4}{7\cdot9}x}{1-\cfrac{\frac{5\cdot8}{9\cdot11}x}{1-\cfrac{\frac{3\cdot6}{11\cdot13}x}{1-\frac{7\cdot10}{13\cdot15}x\ \mathrm{usw.}}}}
\]
so ergibt sich nach Formel 25
\[
\frac1R=G(\frac12,\frac32,\frac72,x),
\qquad\mbox{und}\qquad
\frac1Q=G(\frac52,-\frac12,\frac72,x)
\]
Daraus folgt
\SourceLeaf{147}
\[
\begin{array}{rcl}
RF(\frac12,\frac52,\frac92,x)&=&F(\frac12,\frac32,\frac72,x)\\[4pt]
QF(\frac52,\frac12,\frac92,x)&=&F(\frac52,-\frac12,\frac72,x)
\end{array}
\]
oder, wenn das erste Element mit dem zweiten vertauscht wird,
\[
QF(\frac12,\frac52,\frac92,x)=F(-\frac12,\frac52,\frac72,x)
\]
\par
Nach Gleichung 21 haben wir aber
\[
F(-\frac12,\frac52,\frac72,x)-F(\frac12,\frac32,\frac72,x)
=-\frac47xF(\frac12,\frac52,\frac92,x)
\]
Daraus folgt $Q=R-\frac47x$. Setzt man diesen Wert in die oben angegebene Formel ein, so erhält man
\[
\xi=\frac{\frac2{35}xx}{R-\frac{18}{35}x}
\]
Dies ist die zweite Formel.
\par
Setzt man in Formel 30 $\alpha=\frac mn$, $x=-\gamma nt$, so ergibt sich für einen unendlich großen Wert von $\gamma$
\[
\begin{aligned}
{}[36]\qquad F(\frac mn,1,\gamma,-\gamma nt)
&=1-mt+m(m+n)tt\\&{}-m(m+n)(m+2n)t^3+\mathrm{usw.}\\[8pt]
&=\displaystyle\cfrac{1}{1+\cfrac{mt}{1+\cfrac{nt}{1+\cfrac{(m+n)t}{1+\cfrac{2nt}{1+\cfrac{(m+2n)t}{1+3nt\ \mathrm{usw.}}}}}}}
\end{aligned}
\]
\begin{center}\rule{0.27\textwidth}{0.4pt}\end{center}
